\begin{satz}{Untere Schranke des Erwartungswerts von abs.
\textsl{N}(0,1)--Ord.st.}
{UntereSchrankeDesErwartungswertsVonAbsolutenN01Ordnungsstatistiken}
Seien $n \in \N$, $k \in \haken n$, $\gamma: \Omega \rightarrow \R^n$
unabhängig und $N(0,1)$--verteilt, sowie $u: [0, \infty[\!\ \rightarrow \R, 
t \mapsto \Prim(\abs{\gamma_1}> t)$
(vgl.~\ref{EigenschaftenDerKomplementaerenVerteilungsfunktionVonN01}).\\
Dann ist $\gamma_k^\ast$ nach
\refclaim{EigenschaftenVonOrdnungsstatistiken}
{EigenschaftenVonOrdnungsstatistiken1} messbar, $u$ nach
\refclaim{EigenschaftenDerKomplementaerenVerteilungsfunktionVonN01}
{EigenschaftenDerKomplementaerenVerteilungsfunktionVonN011} invertierbar und es
gilt:
\begin{enumerate}[(1)]
  \item
   \label{UntereSchrankeDesErwartungswertsVonAbsolutenN01Ordnungsstatistiken1}
	$\E\left(\gamma_k^\ast \right) \ge u\inv\left(\frac k {n+1} \right)$.
	\item
	 \label{UntereSchrankeDesErwartungswertsVonAbsolutenN01Ordnungsstatistiken2}
	 $\E\left(\gamma_1^\ast \right)\ge \sqrt{\log(n)}$.
	 \item
	 \label{UntereSchrankeDesErwartungswertsVonAbsolutenN01Ordnungsstatistiken3}
	 Ist $k \le \frac 2 {13}(n+1)$, so ist
	 $\E\left(\gamma_k^\ast \right) \ge \sqrt{\log\left(1+\frac{n+1}k \right)}$.
\end{enumerate}
\end{satz}
\begin{proof}
(1) Offenbar ist 
$\gamma_k^\ast$ die $k$--te Ordnungsstatistik von 
$\pfeil {\abs\cdot} n \circ \gamma$.\\
Es ist $\Prim(\abs{\gamma_1} \le \cdot) = 1 - \Prim(\abs{\gamma_1} > \cdot)
= 1 - u$, also hat $\abs{\gamma_1}$ die stetige Dichte $-u'$.\\
Es gilt für alle $i \in \N_0$:
\begin{align*}
&\frac{n!}{(n+2i+1)!}\cdot\frac{(n-k+2i+1)!}{(n-k)!}
= \left(\prod_{l=0}^{2i}\frac{1}{n+l+1} \right)
\left(\prod_{r=0}^{2i}(n-k+1+r) \right)\\
=&\, \prod_{l=0}^{2i}\frac{n+l+1-k}{n+l+1}
=\prod_{l=0}^{2i}\left(1- \frac{k}{n+l+1} \right)
\ge \prod_{l=0}^{2i}\left(1- \frac{k}{n+1} \right)\\
=&\, \left(1- \frac{k}{n+1} \right)^{2i+1}\text.\tag {$\ast$}
\end{align*}
Sei $(c_j)_{j\in \N_0} \in [0, \infty[^{\N_0}$ die Folge aus
\refclaim{EigenschaftenDerKomplementaerenVerteilungsfunktionVonN01}
         {EigenschaftenDerKomplementaerenVerteilungsfunktionVonN013} und
$B, \Gamma$ wie in~\ref{ElementareEigenschaftenDerGammafunktion}.\\
Somit gilt wegen \refclaim{EigenschaftenVonOrdnungsstatistiken}
{EigenschaftenVonOrdnungsstatistiken2}:
\begin{align*}
\E(\gamma_k^\ast) &\overset{
\refclaim{EigenschaftenVonOrdnungsstatistiken}
{EigenschaftenVonOrdnungsstatistiken4}}= 
\int_0^\infty t \cdot n \binom{n-1}{n-k}\Prim(\abs{\gamma_1}\le t)^{n-k}
(1 - \Prim(\abs{\gamma_1} \le t))^{k-1}(-u'(t))dt\\
&= n \binom{n-1}{n-k} \int_0^\infty t (1-u(t))^{n-k}u(t)^{k-1}(-u'(t))dt\\
&= n \binom{n-1}{n-k} 
\int_{]0, \infty[} (1-u(t))^{n-k}u(t)^{k-1}u\inv(u(t)) \abs{u'(t)}dt\\
&\overset{\substack{\text{Transf.}\\ \text{mit }u}}=
n \binom{n-1}{n-k} \int_{]0,1[}(1-t)^{n-k}t^{k-1}u\inv(t)dt\\
&\overset{\refclaim{EigenschaftenDerKomplementaerenVerteilungsfunktionVonN01}
         {EigenschaftenDerKomplementaerenVerteilungsfunktionVonN013}}=
n \binom{n-1}{n-k} \int_{]0, 1[} (1-t)^{n-k}t^{k-1}\sum_{j=0}^\infty 
c_j(1-t)^{2j+1} dt\\
&=n \binom{n-1}{n-k}\int_{]0,1[} 
\sum_{j=0}^\infty c_j (1-t)^{n+2j+1-k}t^{k-1}dt\\
&\overset{\substack{\text{Satz v. d.}\\\text{mon. Konv.}}}= 
n \binom{n-1}{n-k} \sum_{j=0}^\infty c_j
\int_{]0,1[}(1-t)^{n+2j+1-k}t^{k-1}dt\\
&=n \binom{n-1}{n-k} \sum_{j=0}^\infty c_j
\int_0^1(1-t)^{n+2j+1-k}t^{k-1}dt\\
&\overset{\text{Def. }B}= n \binom{n-1}{n-k} \sum_{j=0}^\infty c_j
B(n+2j-k+2, k)\\
&\overset{\refclaim{ElementareEigenschaftenDerGammafunktion}
{ElementareEigenschaftenDerGammafunktion4}}=
n \binom{n-1}{n-k} \sum_{j=0}^\infty c_j \frac{\Gamma(n+2j-k+2)\Gamma(k)}
{\Gamma(n+2j+2)}\\
&\overset{\refclaim{ElementareEigenschaftenDerGammafunktion}
{ElementareEigenschaftenDerGammafunktion3}}= 
\sum_{j=0}^\infty c_j \frac{n!}{(k-1)!(n-k)!}
\frac{(k-1)!(n-k+2j+1)!}{(n+2j+1)!}\\
&= \sum_{j=0}^\infty c_j \frac{n!}{(n-k)!}\frac{(n-k+2j+1)!}{(n+2j+1)!}
\overset{(\ast)}\ge \sum_{j=0}^\infty c_j
\left(1-\frac{k}{n+1}\right)^{2j+1}\\
&\overset{\refclaim{EigenschaftenDerKomplementaerenVerteilungsfunktionVonN01}
{EigenschaftenDerKomplementaerenVerteilungsfunktionVonN013}}= 
u\inv\left(\frac k {n+1} \right)\text.
\end{align*}
(2) Es gilt:
$\E(\gamma_1^\ast) \overset{(1)}\ge u\inv\left(\frac 1 {n+1} \right)
\overset{\refclaim{EigenschaftenDerKomplementaerenVerteilungsfunktionVonN01}
{EigenschaftenDerKomplementaerenVerteilungsfunktionVonN015}}\ge
\sqrt{\log(n)}$.\\
(3) Es gelte $k \le \frac 2{13}(n+1)$.\\
Dann ist $\frac{k}{n+1} \le \frac 2 {13}$, also
$\E(\gamma_k^\ast) \overset{(1)}\ge u\inv\left(\frac k {n+1} \right)
\overset{\refclaim{EigenschaftenDerKomplementaerenVerteilungsfunktionVonN01}
{EigenschaftenDerKomplementaerenVerteilungsfunktionVonN016}}\ge
\sqrt{\log\left(1 + \frac{n+1}k \right)}$.
\end{proof}